\documentclass[10pt,a4paper]{amsart}
\usepackage[margin=19mm]{geometry}
\usepackage{amsmath,amssymb,mathtools}
\usepackage[scr=rsfs,cal=euler]{mathalfa}
\usepackage{xfrac}
\usepackage[unicode,hidelinks,bookmarks=false]{hyperref}
\newcommand{\ie}{i.e.\ }
\newcommand{\cf}{cf.\ }
\newcommand{\resp}{respectively\ }
\newcommand{\Subsection}{\subsection}
\providecommand{\nobreakdash}{\nobreak\hbox{-}}
\setlength{\emergencystretch}{2em}
\multlinegap0pt
\usepackage{bm}
\usepackage{bbm}
\usepackage{dsfont}
\usepackage{xspace}
\usepackage{cancel}
\usepackage{mathtools}
\usepackage{amsthm}
\makeatletter
\@ifundefined{theorem}{\newtheorem{theorem}{Theorem}}{}
\@ifundefined{lemma}{\newtheorem{lemma}[theorem]{Lemma}}{}
\makeatother
\def\ker{\textrm{ker}}
\def\res{\mathrm{res}}
\newcommand{\C}{\mathbb C}
\newcommand{\End}{\mathcal E nd}
\newcommand{\id}{\mathrm{id}}
\newcommand{\tr}{\mathrm{tr}}
\title[Symplectic leaves of meromorphic Hitchin systems]{Symplectic leaves of meromorphic Hitchin systems}
\author{Jia Choon Lee, Sukjoo Lee}
\date{Source: arXiv:2606.31123v1 (CC BY 4.0)}
\begin{document}
\maketitle

\noindent\emph{Source and licence.} Symplectic leaves of meromorphic Hitchin systems. Version arXiv:2606.31123v1 (CC BY 4.0), distributed under the Creative Commons Attribution 4.0 International licence, as recorded in the arXiv licence metadata for this exact version and shown on the abstract page https://arxiv.org/abs/2606.31123v1. Full rights notice: licenses/arxiv-2606.31123-CC-BY-4.0.txt.\par\smallskip

\noindent\textit{Editorial scope.} Counted unit: the Lemma whose source label is \texttt{lem:markman-residue-submersion} in the article (source0.tex lines 745--747, source label \texttt{lem:markman-residue-submersion}), delivered together with the article's own complete original proof of it (source0.tex lines 749--806, one \texttt{proof} environment of 58 lines). No mathematics is written, repaired or re-derived: statement and proof are the article's own text line for line. Required context retained uncounted: none. Every definition, display and label of those lines is retained unchanged. Omitted from this package and not redistributed: 1--744, 748--748, 807--2148. Nothing inside a retained excerpt is omitted or altered: every retained body is delivered exactly as the article's lines are written, so no omission receipt of any kind accompanies this package. Citations: the retained excerpts carry no citation command at all, so no bibliography entry of the article is transcribed and no reference is redistributed. Reference adaptation: none was needed and none was made; every reference inside the delivered unit resolves inside it, so no unresolved reference or citation is produced. Printed slips kept literal: no printed word, symbol, hypothesis or number of the counted statement or of its proof was altered. Layout and class: the article's own class is \texttt{amsart} with packages including inputenc, lmodern, amsmath, amsthm, amssymb, amsfonts, bm, ulem, hyperref, varwidth, adjustbox, enumerate, verbatim, longtable; the wrapper is the lane's pinned \texttt{amsart} editorial wrapper at 10pt on A4, which restates the theorem-like environments the retained text needs and adds the article's own macro definitions that the retained text needs (\texttt{\textbackslash C}, \texttt{\textbackslash End}, \texttt{\textbackslash id}, \texttt{\textbackslash ker}, \texttt{\textbackslash res}, \texttt{\textbackslash tr}), with extra wrapper packages (bm, bbm, dsfont, xspace, cancel, mathtools). The delivered numbering is the unit's own. Assets: no retained span contains a figure environment, a graphics-inclusion command, a PDF-page inclusion, a table or a TikZ picture; no image file is copied into this package, the package declares no asset dependency and no image file is redistributed with it. Licence: version arXiv:2606.31123v1 of this article is distributed under the Creative Commons Attribution 4.0 International licence, as recorded in the arXiv licence metadata for this exact version and shown on the abstract page https://arxiv.org/abs/2606.31123v1. A full rights notice is delivered at licenses/arxiv-2606.31123-CC-BY-4.0.txt. Characters: every retained excerpt is pure ASCII, so the delivered engine, XeLaTeX with its default Latin Modern text font, reports no missing character.
\begin{lemma}\label{lem:markman-residue-submersion}
	The residue map $\mu:T^*\mathcal U_D\to\mathfrak R$ is smooth.
\end{lemma}

\begin{proof}
	It is enough to show that $d\mu$ is surjective at every point. Let
	$y=(E,\Phi,\eta)$. Since $E$ is stable, $H^0(C,\End(E))=\C\cdot\id_E$. Consider the residue map for Higgs-field variations
	\[
	\res_D:
	H^0(C,\End(E)\otimes K_C(D))
	\longrightarrow
	\bigoplus_i\End(E_{p_i}).
	\]
	Its image is contained in the  residue space
	\[
	\mathfrak R_E:=
	\left\{
	(B_i)_i\in\bigoplus_i\End(E_{p_i})
	\ \middle|\
	\sum_i\tr(B_i)=0
	\right\}.
	\]
	We claim that $\res_D$ is surjective onto $\mathfrak R_E$. Indeed, the residue sequence
	\[
	0\to\End(E)\otimes K_C
	\to\End(E)\otimes K_C(D)
	\to \bigoplus_i\End(E_{p_i}) \to 0
	\]
	gives a connecting homomorphism $\delta:\bigoplus_i\End(E_{p_i})
	\longrightarrow H^1(C,\End(E)\otimes K_C)$,
	and $\operatorname{coker}(\res_D)$ identifies with $\operatorname{Im}(\delta)$.
	By Serre duality and the trace pairing, the dual of $\delta$ is the evaluation
	map
	\[
	H^0(C,\End(E))
	\longrightarrow
	\left(\bigoplus_i\End(E_{p_i})\right)^*,
	\qquad
	s\longmapsto \bigl((B_i)_i\mapsto \sum_i\tr(s(p_i)B_i)\bigr).
	\]
    where $(-)^*$ denotes the linear dual. By the assumption that $E$ is stable, $H^0(C,\End(E))=\C\cdot\id_E$, hence the annihilator of
	$\operatorname{Im}(\res_D)=\ker(\delta)$ is spanned by the functional $(B_i)_i\longmapsto \sum_i\tr(B_i)$. Thus
	\[
	\operatorname{Im}(\res_D)
	=
	\left\{
	(B_i)_i\in\bigoplus_i\End(E_{p_i})
	\ \middle|\
	\sum_i\tr(B_i)=0
	\right\}
	=\mathfrak R_E.
	\]

	Now vary only the Higgs field, keeping $E$ and the framing $\eta$ fixed:
	$\Phi_\varepsilon=\Phi+\varepsilon\psi$, where
	$\psi\in H^0(C,\End(E)\otimes K_C(D))$. Then $
	d\mu_y(\psi)
	=
	\bigl(\eta_i\circ\res_{p_i}(\psi)\circ\eta_i^{-1}\bigr)_i$. Conjugation by the framings $\eta_i$ identifies $\mathfrak R_E$ with
	$\mathfrak R$. Since $\res_D$ is surjective onto $\mathfrak R_E$, the image of
	$d\mu_y$ contains all of $\mathfrak R$. Hence $d\mu_y$ is surjective.
\end{proof}

\end{document}
